% Complete Chapter 5 introduction, source PDF 61 / printed 51.
% block: approximation-hilbert-heading page:61 type:chapter
\clearpage\setcounter{section}{4}
\section{\BT{Approximation of Hilbert complexes}{Hilbert 复形的逼近}}\label{chap:approximation-hilbert-complexes}

% block: approximation-hilbert-goal page:61 type:ordinary-paragraph
\BT{Our goal now is to design and analyze numerical methods for the approximate solution of a Hodge Laplacian problem associated to some closed Hilbert complex. We shall use Galerkin methods, i.e., we shall replace the trial and test spaces occuring in the primal or mixed weak formulation by finite dimensional subspaces. In practice, when the Hilbert complex is the de Rham complex or something similar, the subspaces we consider will be finite element spaces.}{现在的目标是设计和分析数值方法, 以近似求解与某个闭 Hilbert 复形相伴的 Hodge Laplacian 问题. 我们将采用 Galerkin 方法, 即用有限维子空间替换原始或混合弱表述中出现的试探空间和检验空间. 在实际应用中, 当 Hilbert 复形为 de Rham 复形或类似复形时, 所考虑的子空间将是有限元空间.}

% block: approximation-short-complex page:61 type:ordinary-paragraph
\BT{To define the Hodge Laplace problem at some level $k$ of the Hilbert complex, we require only the operators $d^{k-1}$ from $W^{k-1}$ to $W^k$ with domain $V^{k-1}$ and the $d^k$ from $W^k$ to $W^{k+1}$ with domain $V^k$; cf. \eqref{eq:abstract-hodge-operator}. Therefore, in the interest of carrying no more baggage than necessary, in this chapter we do not require the whole Hilbert complex \eqref{eq:hilbert-domain-complex} but only a short segment of it:}{要定义 Hilbert 复形在某个层次 $k$ 上的 Hodge Laplace 问题, 只需要从 $W^{k-1}$ 到 $W^k$, 定义域为 $V^{k-1}$ 的算子 $d^{k-1}$, 以及从 $W^k$ 到 $W^{k+1}$, 定义域为 $V^k$ 的算子 $d^k$; 参见 \eqref{eq:abstract-hodge-operator}. 因此, 为避免引入不必要的负担, 本章不需要整个 Hilbert 复形 \eqref{eq:hilbert-domain-complex}, 而只需要其中的一个短片段:}
\begin{equation}\label{eq:approximation-short-complex}
V^{k-1}\xrightarrow{d}V^k\xrightarrow{d}V^{k+1}.
\end{equation}

% block: approximation-chapter-outline page:61 type:ordinary-paragraph
\BT{We shall see that Galerkin discretization of the primal formulation is problematic and shall focus on Galerkin discretization of the mixed formulation. This requires two compatible Galerkin subspaces, one for $V^{k-1}$ and one for $V^k$. The main result of this chapter is to establish assumptions on the choice of subspaces, which ensures consistency and stability of the Galerkin method, and to deduce consequences from it. Box~\ref{box:approximation-hilbert-summary} summarizes the main results of the chapter.}{我们将看到, 原始表述的 Galerkin 离散化存在问题, 因此将着重讨论混合表述的 Galerkin 离散化. 这需要两个相容的 Galerkin 子空间, 分别对应 $V^{k-1}$ 和 $V^k$. 本章的主要结果是对子空间的选取建立假设, 保证 Galerkin 方法的相容性和稳定性, 并由此导出推论. Box~\ref{box:approximation-hilbert-summary} 概括了本章的主要结果.}
